% ============================================================
% templates/t05-cite.tex   —— T5 作品 / 论文条目
%
% 适用场景：一条记录占两行 —— 标题行（可带链接）+ 说明行（角色 / 出处）。
%          论文列表、项目经历、作品集、获奖记录。
%          与 T1 的区别：T1 是一行一条，本模板的"一条"有两行，
%          而且两行要读起来是一个整体。
% 提供    ：\DeckCite
% 依赖    ：\DeckLink（primitives.tex）
% 参数    ：#1 标题行（可内嵌 \DeckLink{URL}{文字}）
%           #2 说明行（角色 / 出处 / 年份）
% 容量    ：一页最多 5 条（每条两行 + 条间距，比 T1 少）
%           标题行一行 = \SlideLineUnitsSafe 个全角字（同 T1）
%           说明行同理，但它通常更短，不必压到极限
% 易错    ：① 两行之间 0.12cm、条目之间 0.34cm —— 这个差是刻意的，
%             让两行读作一条、条目之间有明确呼吸。不要自己插 \DeckGap
%           ② 链接不要写成裸 URL 文本，用 \DeckLink 包起来才有强调色
%           ③ 编译通过不代表链接生效 —— 要验证请用 pdftohtml 查 <a href=，
%             别在 PDF 原始字节里搜（对象流压缩会让它搜不到）
%
% 【填满的样子】
%   \DeckCite
%     {基于模板库的幻灯片排版方法}
%     {第一作者，2024}
%
%   \DeckCite
%     {\DeckLink{https://example.com}{项目主页}}
%     {开源项目，2023}
% ============================================================

% --- 条目：一条两行（标题行 + 说明行）---
% #1 标题行（可内嵌 \DeckLink），#2 说明行（角色 / 出处）
% 两行之间 0.12cm，让它们读作一条完整记录；条目之间 0.34cm。
\DeckDefine{\DeckCite}[2]{%
  \par
  \hangindent=1.1em\hangafter=1
  \noindent #1\par
  \vspace*{0.12cm}%
  \noindent #2\par
  \vspace*{0.34cm}%
}
